\documentclass[a4paper,twoside]{article}

\usepackage{morewrites}

\usepackage[svgnames,dvipsnames,x11names]{xcolor}
\usepackage{graphicx}
\usepackage[english]{babel}
\usepackage[colorlinks=true, linkcolor=NavyBlue, urlcolor=NavyBlue, citecolor=NavyBlue]{hyperref}
\usepackage[a4paper]{geometry}
\usepackage{ts-fonts}
\usepackage{amsmath}
\usepackage{float}
\usepackage{seqsplit}
\usepackage{ts-code}

\usepackage{lastpage}
\usepackage{refcount}
\makeatletter
\def\ps@plain{
  \let\@oddhead\@empty
  \let\@evenhead\@empty
  \def\@oddfoot{
    \hfil
    \ifnum\getpagerefnumber{LastPage}>1\relax
      \thepage
    \fi
    \hfil}
  \let\@evenfoot\@oddfoot
}
\makeatother

\title{Images}
\author{}\date{}

\begin{document}

\maketitle

\thispagestyle{plain}
Images can be included in Markdown using the following syntax:

\begin{code}{md}{}{}
\begin{Verbatim}[commandchars=\\\{\},breaklines, breakanywhere, commandchars=\\\{\}]
![\PY{n+nt}{Alt text}](\PY{n+na}{https://picsum.photos/400/150})\PYZob{}width=50\PYZpc{}\PYZcb{}
\end{Verbatim}
\end{code}
The width attribute is useful to scale images directly in the Markdown source.

\begin{figure}[H]
    \centering

    \ifdefined\adjustbox

        \adjustbox{max width=\textwidth}{
            \includegraphics[width=\linewidth]{snippet-<HASH>.pdf}
        }

    \else

        \includegraphics[width=\linewidth]{snippet-<HASH>.pdf}

    \fi

\end{figure}
\section{Draw.io diagrams}\label{draw-io-diagrams}

A \texttt{.drawio} file included as an image is exported to a vector PDF during the
build:

\begin{code}{md}{}{}
\begin{Verbatim}[commandchars=\\\{\},breaklines, breakanywhere, commandchars=\\\{\}]
![\PY{n+nt}{Euclidean GCD}](\PY{n+na}{pgcd.drawio})\PYZob{}width=60\PYZpc{}\PYZcb{}
\end{Verbatim}
\end{code}
The export is \emph{cropped} to the drawing by default --- the surrounding canvas the
diagram sits on is discarded, so the figure carries no unexpected white margin
and \texttt{width=} scales the drawing itself. When the page layout is part of the
diagram (a frame, a title block, deliberate margins, several drawings placed at
fixed positions on the same sheet), ask for the page instead with \texttt{crop=false}:

\begin{code}{md}{}{}
\begin{Verbatim}[commandchars=\\\{\},breaklines, breakanywhere, commandchars=\\\{\}]
![\PY{n+nt}{Site plan}](\PY{n+na}{plan.drawio})\PYZob{}width=100\PYZpc{} crop=false\PYZcb{}
\end{Verbatim}
\end{code}
This is draw.io's \emph{Size: Page Size} export: the output covers the page(s) the
drawing spans, at the \texttt{pageWidth} / \texttt{pageHeight} of the diagram, whichever
backend (\texttt{playwright}, \texttt{local}, \texttt{docker}) performs the conversion.

The same diagram can appear both ways in one document; each variant is exported
and stored as its own asset.

To flip the default for a whole document, set \texttt{drawio\_crop} in the front matter
(or \texttt{-\allowbreak{}-\allowbreak{}attribute press.drawio\_crop=false} on the command line); an image
attribute still wins over it:

\begin{code}{yaml}{}{}
\begin{Verbatim}[commandchars=\\\{\},breaklines, breakanywhere, commandchars=\\\{\}]
\PY{n+nn}{\PYZhy{}\PYZhy{}\PYZhy{}}
\PY{n+nt}{press}\PY{p}{:}
\PY{+w}{  }\PY{n+nt}{drawio\PYZus{}crop}\PY{p}{:}\PY{+w}{ }\PY{l+lScalar+lScalarPlain}{false}
\PY{n+nn}{\PYZhy{}\PYZhy{}\PYZhy{}}
\end{Verbatim}
\end{code}

\end{document}
